\section{Curvature estimates and the Pfaffian bound}
\label{sec:curvature}

We estimate the curvature of the connection constructed in
Proposition~\ref{co:connection}.  The smooth stages are controlled by
operators on weighted Hilbert spaces.  The raw stage is instead controlled
by the number of nearby simplex coordinates.  These two estimates will be
combined only after averaging over the coordinate subsets in the
Pfaffian expansion.

Throughout this section, the word metric and the fixed construction
constants are understood.  The constant $C$ may depend on these data but
not on $R,n,g$, or the tuple.  A constant $C_R$ may additionally depend
on $R$.  All parameters $r,W,b_*,\beta_j,\delta_j$, and $a_j$ retain
their definitions from Section~\ref{sec:connection}.  In particular,
$\delta_j=\beta_j/n$ and $10\delta_{b_*}<1$.

We work in the section $t_0=0$, on the simplex
\[
 \Delta_n=\{0\le t_1\le\cdots\le t_n\le1\},
 \qquad n=2m.
\]
Expectation with respect to uniform simplex probability will be denoted
by $\mathbb E_{\Delta_n}$; its density in these coordinates is $n!$.
If $Q$ is a two-form, its matrix is defined by
$Q=\sum_{i<i'}Q_{ii'}\,dt_i\wedge dt_{i'}$.  Our convention is
\begin{equation}
 \frac{Q^m}{m!}
 =\Pf(Q)\,dt_1\wedge\cdots\wedge dt_n.
 \label{pf:normalization}
\end{equation}
Equation~\eqref{cy:factorials} identifies the expectation below with
the normalized integral in Proposition~\ref{cy:trace-identity}.  For a subset
$I\subset\{1,\ldots,n\}$ of even cardinality, $\Pf_I(Q)$ means the
Pfaffian of the principal matrix on $I$, listed in increasing order.
We set $\Pf_{\varnothing}(Q)=1$.

\begin{theorem}[Averaged curvature bound]
\label{pf:bound}
Fix a sufficiently large $R$, and set $S=b_*+2$.  Let $n$ be an even
integer with $n>10\beta_{b_*}$.  Suppose either that $g$ has infinite
order, or that $g$ has finite order and
\[
 \min_{h\in G}d(h,gh)>4r.
\]
For every tuple $(h_0,\ldots,h_n)$, the connection $\alpha=B/C_{\rm nor}$
of Proposition~\ref{co:connection} satisfies
\begin{equation}
 \mathbb E_{\Delta_n}|\Pf(d\alpha)|
 \le C_R\bigl(CS^{3/2}R^{-1/10}\bigr)^n
          \prod_{i=1}^n l_i,
 \qquad l_i=1+d(h_{i-1},h_i).
 \label{pf:main-estimate}
\end{equation}
The constants are independent of $n,g$, and the tuple.  Moreover,
$S=O(1+\log R)$ as $R\to\infty$.
\end{theorem}

We prove the theorem in three steps: estimates for principal Pfaffians
of the smooth stages, an averaged estimate for the raw stage, and the
expansion of the normalized curvature.

\subsection{Smooth stages as weighted operators}

Fix a stage $j>0$ and a point of the simplex.  Let $\Lambda_j$ be its
label set and equip $(0,1)\times\Lambda_j$ with the measure
\[
 d\vartheta_j(u,y)=Y_j(u)_y\,du.
\]
We use the real Hilbert space $\mathcal H_j=L^2(\vartheta_j)$.  Recall that
the total mass of $Y_j(u)$ is at most two, and is at most one at the
last stage.  Define the operator $T_j$ by the kernel
\begin{equation}
 T_j(u,y;v,z)
   =\sum_{a\in\Z}U_j(v+a-u)K_j(y,g^a z),
 \label{pf:kernel}
\end{equation}
where the action of $g$ on labels is the one used in the channel
construction.  Oddness of $U_j$, symmetry of $K_j$, and simultaneous
equivariance of $K_j$ show that this kernel is skew.  Indeed, exchange
$(u,y)$ with $(v,z)$ and replace $a$ by $-a$ in the sum.

For $1\le i\le n$, put
\[
 X_i(u,y)=\frac{\partial_iY_j(u)_y}{Y_j(u)_y}
 \quad\text{when }Y_j(u)_y>0,
\]
and give $X_i$ the value zero at the remaining points.  The derivatives
of the zero weights vanish, and \eqref{co:fisher} gives
\begin{equation}
 \|X_i\|_{\mathcal H_j}
 \le \frac{a_jl_i}{\sqrt{\delta_j}}.
 \label{pf:vector-norm}
\end{equation}
The exterior-derivative identity \eqref{co:exterior-pairing} is
\[
 dB_j=\frac12\int_{\rm per}du\int_{\R}dv\,
       U_j(v-u)K_j\bigl(dY_j(u)\wedge dY_j(v)\bigr).
\]
The integral over a period can be taken with fixed endpoints because
its integrand is periodic.  Reducing $v$ to $(0,1)$ gives
\eqref{pf:kernel}; the two antisymmetric terms then agree.  Consequently,
\begin{equation}
 (dB_j)_{ii'}=\langle X_i,T_jX_{i'}\rangle_{\mathcal H_j}.
 \label{pf:curvature-pairing}
\end{equation}

Here is the finite-dimensional estimate we shall apply to this pairing.
Let $I$ have even size $k>0$, let $E$ be the span of the vectors $X_i$,
$i\in I$, and let $s_1\ge s_2\ge\cdots$ be the singular values of the
compression $P_ET_j|_E$, completed by zeros.  Then
\begin{equation}
 |\Pf_I(dB_j)|
 \le \prod_{i\in I}\|X_i\|
       \left(\prod_{\nu=1}^k s_\nu\right)^{1/2}.
 \label{pf:compression}
\end{equation}
To see this, the matrix in \eqref{pf:curvature-pairing} has rank less
than $k$ if $\dim E<k$, so its Pfaffian is zero.  Otherwise write the
$X_i$ as columns of a square matrix $V$ in an orthonormal basis of $E$.
The curvature matrix is $V^{\mathsf T}T_EV$.  Taking determinants and
using $\Pf(A)^2=\det A$ for a skew matrix gives the right side of
\eqref{pf:compression}, after Hadamard's inequality
$|\det V|\le\prod_i\|X_i\|$.

For $1\le j<b_*$, the support of $U_j$ is contained in
$[-2\delta_{j+1},2\delta_{j+1}]$.  At most one integer $a$ in
\eqref{pf:kernel} contributes for a fixed pair $(u,v)$, because
$4\delta_{j+1}<1$.  Thus the kernel is bounded by a fixed constant and
vanishes unless the circular distance between $u$ and $v$ is at most
$2\delta_{j+1}$.  The mass bound on the channels implies
\[
 \|T_j\|_{\rm HS}^2
 =\int\!\int |T_j(u,y;v,z)|^2\,
       d\vartheta_j(u,y)d\vartheta_j(v,z)
 \le C\delta_{j+1}.
\]
Bessel's inequality gives the same upper bound for the squared
Hilbert--Schmidt norm of every compression.  Applying the arithmetic
mean inequality to $s_1^2,\ldots,s_k^2$ gives
\[
 \left(\prod_{\nu=1}^k s_\nu\right)^{1/2}
 \le C^k\delta_{j+1}^{k/4}k^{-k/4}.
\]
Combining this with \eqref{pf:vector-norm} and
\eqref{pf:compression}, and using $\beta_{j+1}=\beta_j^{3/2}$, proves
\begin{equation}
 |\Pf_I(dB_j)|
 \le\bigl(Ca_j\beta_j^{-1/8}\bigr)^k
          (n/k)^{k/4}\prod_{i\in I}l_i,
 \qquad 1\le j<b_*.
 \label{pf:intermediate}
\end{equation}
In fact, the powers of the time scales simplify as
\[
 \delta_j^{-k/2}\delta_{j+1}^{k/4}k^{-k/4}
 =\beta_j^{-k/8}(n/k)^{k/4}.
\]

\subsection{The terminal sign kernel}

The last lag kernel is not supported at a small time scale.  Its useful
property is instead that an order comparison is constant between two
disjoint ordered intervals.  The following elementary approximation
works even when the source and target measures differ.

\begin{lemma}[Approximation of order kernels]
\label{pf:order-kernel}
Let $\mu,\nu$ be finite nonatomic positive measures on $\R$, and set
$M=\mu(\R)+\nu(\R)$.  Let $A:L^2(\nu)\to L^2(\mu)$ have any one of
the kernels
\[
 \operatorname{sign}(s-t),\qquad
 \mathbf1_{s<t},\qquad \mathbf1_{s>t},
\]
where $t$ is the output variable and $s$ the input variable.  For each
integer $N\ge1$, there is an operator $A_N$ of rank at most $N$ such
that
\[
 \|A-A_N\|\le \frac{M}{2N}.
\]
Also $\|A\|\le\sqrt{\mu(\R)\nu(\R)}\le M/2$.
\end{lemma}

\begin{proof}
If $M=0$, the assertion is immediate.  Otherwise partition $\R$ into
$N$ successive intervals $I_1,\ldots,I_N$, each having
$(\mu+\nu)$-measure $M/N$.  Such a partition exists because the measure
is nonatomic.  Endpoint conventions do not affect the operators.
Let $A_N$ retain the kernel of $A$ when $t$ and $s$ lie in different
intervals, and set the kernel to zero when they lie in the same
interval.  On $I_a\times I_b$ with $a\ne b$, the order of $t$ and $s$
is fixed.  The range of $A_N$ is therefore contained in the span of the
$N$ functions $\mathbf1_{I_a}$ in $L^2(\mu)$.

The remainder $A-A_N$ maps $L^2(\nu|_{I_a})$ to
$L^2(\mu|_{I_a})$ for each $a$, without mixing different intervals.
On that block its kernel is bounded by one, so Cauchy--Schwarz gives
the operator norm bound
\[
 \sqrt{\mu(I_a)\nu(I_a)}
 \le \frac{\mu(I_a)+\nu(I_a)}2=\frac{M}{2N}.
\]
Both the source blocks and the target blocks are orthogonal.  Hence
the norm of their direct sum is the maximum of their norms, proving
the approximation claim.  Applying the same Cauchy--Schwarz estimate
on the whole line gives the asserted bound on $\|A\|$.
\end{proof}

\begin{lemma}[Singular values at the last stage]
\label{pf:terminal-singular}
For the operator $T_{b_*}$, every finite-dimensional orthogonal
compression has singular values satisfying
\begin{equation}
 s_\nu\le \frac C\nu,\qquad \nu\ge1.
 \label{pf:singular-decay}
\end{equation}
The constant is independent of $n,R,g$, and the tuple, in either of
the two cases of Theorem~\ref{pf:bound}.
\end{lemma}

\begin{proof}
We will approximate $T_{b_*}$, for every $x>0$, by an operator of rank
at most $C/x$ and norm error at most $x$.

First suppose that $g$ has infinite order.  The final label set is $G$
and the label kernel is equality.  It follows that
$\mathcal H_{b_*}$ splits orthogonally according to the left
$g$-orbits in $G$.  In one orbit, choose a representative $y_0$ and
write a label as $g^a y_0$, $a\in\Z$.  Associate to the pair
$(u,g^a y_0)$ the real coordinate $u-a$.  If the second label is
$g^b y_0$, the equality kernel in \eqref{pf:kernel} selects precisely
the shift $a-b$.  The resulting kernel is
\[
 \operatorname{sign}(v+a-b-u)
 =\operatorname{sign}\bigl((v-b)-(u-a)\bigr).
\]
Thus this orbit gives a sign operator on a nonatomic measure on the
line.  Let $b$ be the mass of that measure.  The orbit masses sum to
at most one, because the final channel is a subprobability channel.
Lemma~\ref{pf:order-kernel}, with the same measure in source and
target, gives norm at most $b$ and a rank-$N$ approximation with error
at most $b/N$.

Discard each orbit block with $b\le x$.  On every other block choose
$N=\lceil b/x\rceil$.  The error of the resulting direct-sum
approximation is at most $x$.  Moreover,
\[
 \sum_{b>x}\lceil b/x\rceil
 \le \frac1x\sum_{b>x}b+\#\{b>x\}
 \le \frac2x.
\]
This proves the desired rank bound in the infinite-order case.

Now suppose that $g$ has finite order.  Put
$d\mu_y(u)=Y_{b_*}(u)_y\,du$ on $(0,1)$, extend this measure by zero
to $\R$, and write $\mathcal H_y=L^2(\mu_y)$.
Since the final lag kernel is
$\operatorname{sign}(s)\mathbf1_{|s|<1}$, the only possible shifts for
$0<u,v<1$ are $0,1,-1$.  The operator is the sum of the three kernels
\begin{equation}
 \begin{split}
  &\operatorname{sign}(v-u)\mathbf1_{y=z},\\
  &\mathbf1_{v<u}\mathbf1_{y=gz},\\
  &-\mathbf1_{v>u}\mathbf1_{y=g^{-1}z}.
 \end{split}
 \label{pf:three-shifts}
\end{equation}
The shifted contributions are direct sums of maps
$\mathcal H_z\to\mathcal H_{gz}$ and
$\mathcal H_z\to\mathcal H_{g^{-1}z}$, respectively.  After reindexing
the target summands, each is an orthogonal direct sum: its source
blocks are mutually orthogonal and its target blocks are mutually
orthogonal.  The identity contribution has the same property.
This remains true even though the measures in a paired block need
not coincide; the shifted summands need not be invariant subspaces.

For a paired block, let $b=\mu_y(\R)+\mu_z(\R)$.  By
Lemma~\ref{pf:order-kernel}, its norm is at most $b$, and it has a
rank-$N$ approximation with error at most $b/N$.  The sums of these
$b$'s over a single contribution are at most two: every label mass
appears once as a source mass and once as a target mass.  Discarding
blocks of mass at most $\varepsilon$ and taking
$N=\lceil b/\varepsilon\rceil$ on the others therefore gives error at
most $\varepsilon$ and total rank at most $4/\varepsilon$.  Use
$\varepsilon=x/3$ for each contribution in \eqref{pf:three-shifts}
and add the three approximations.  Their total error is at most $x$
and their total rank is at most $C/x$.  If $g$ has order two, the two
shift permutations coincide; the separate estimates still apply and
are added in exactly the same way.

Finally, orthogonal compression cannot increase either the rank or
the error of an approximation.  Singular values beyond the dimension
of the compressed space are zero, so consider an index $\nu$ no larger
than that dimension.  If the compressed matrix has an approximation
of rank less than $\nu$ with error $x$, its $\nu$th singular value is
at most $x$.  Indeed, the span of its first $\nu$
right singular vectors contains a unit vector annihilated by the
approximating matrix, and the original matrix has norm at least
$s_\nu$ on that vector.  If the rank bound above is $C_0/x$, take
$x=2C_0/\nu$.  The approximating rank is then at most $\nu/2<\nu$,
which proves \eqref{pf:singular-decay}.
\end{proof}

For an even principal set $I$ of size $k>0$, insert
\eqref{pf:singular-decay} in \eqref{pf:compression}.  Multiplication
of the two-form by $W$ multiplies its size-$k$ Pfaffian by $W^{k/2}$.
We obtain
\[
 |\Pf_I(WdB_{b_*})|
 \le W^{k/2}\left(\frac{a_{b_*}}{\sqrt{\delta_{b_*}}}\right)^k
       \frac{C^{k/2}}{\sqrt{k!}}\prod_{i\in I}l_i.
\]
For $k\ge1$, integration of $\log x$ gives
$\log(k!)\ge\int_1^k\log x\,dx\ge k\log k-k$, and hence
$k!\ge(k/e)^k$.  We conclude that
\begin{equation}
 |\Pf_I(WdB_{b_*})|
 \le\bigl(Ca_{b_*}\sqrt W\,\beta_{b_*}^{-1/2}\bigr)^k
          (n/k)^{k/2}\prod_{i\in I}l_i.
 \label{pf:terminal}
\end{equation}
The estimates \eqref{pf:intermediate} and \eqref{pf:terminal} control
every smooth stage at every point of the simplex.  It remains to
control the raw stage in the average required for the full Pfaffian.

\subsection{The raw stage and nearby coordinates}

Recall $J_i=\delta_{h_{i-1}}-\delta_{h_i}$.  The coefficient of $dt_i$
in $B_0$ is
\[
 -\frac12\int_{\R}U_0(v-t_i)K_0(J_i,Y_0(v))\,dv.
\]
For distinct $i,i'$ in the interior of the simplex, differentiating
at the boundary $t_{i'}$ and antisymmetrizing gives
\[
 (dB_0)_{ii'}
 =\sum_{a\in\Z}U_0(t_{i'}+a-t_i)K_0(J_i,g^aJ_{i'}).
\]
There is at most one nonzero summand.  Since
$K_0(h,h')=\chi_R(d(h,h'))$ and $\chi_R$ is $1/R$-Lipschitz,
\[
 |K_0(J_i,g^aJ_{i'})|
 \le\frac{2d(h_{i-1},h_i)}R
 \le\frac{2l_i l_{i'}}R.
\]
The support of $U_0$ also shows that this matrix entry is zero if the
circular distance between $t_i$ and $t_{i'}$ exceeds $2R/n$.

For $I\subset\{1,\ldots,n\}$ and $i\in I$, let $N_i(I)$ be the
number of indices in $I$, including $i$ itself, whose times have
circular distance at most $2R/n$ from $t_i$.  If $k=|I|$ is even,
divide row and column $i$ of the principal curvature matrix by $l_i$.
Each resulting row has Euclidean norm at most
$(C/R)\sqrt{N_i(I)}$.  Hadamard's inequality, followed by the
Pfaffian--determinant identity, yields
\begin{equation}
 |\Pf_I(dB_0)|
 \le (C/R)^{k/2}\prod_{i\in I}l_i
                    \prod_{i\in I}N_i(I)^{1/4}.
 \label{pf:raw}
\end{equation}

We next average the last product.  The needed assertion is an average
over subsets of a fixed size; it is not a bound asserted separately
for each fixed subset of the ordered coordinates.

\begin{lemma}[Average of the neighbor counts]
\label{pf:occupancy}
Let $R\ge1$, $n\ge8R$, and $0\le k\le n$.  Choose
$(t_1,\ldots,t_n)$ with uniform probability on $\Delta_n$, and
independently choose $I$ uniformly among the $k$-element subsets of
$\{1,\ldots,n\}$.  Then
\begin{equation}
 \mathbb E_{\Delta_n,I}\prod_{i\in I}N_i(I)^{1/4}
 \le C^k R^{k/4},
 \label{pf:count-average}
\end{equation}
with the empty product interpreted as one.
\end{lemma}

\begin{proof}
The statement is immediate for $k=0$.  Uniform simplex coordinates
are the order statistics of $n$ independent uniform points in
$(0,1)$.  Selecting a uniform set of $k$ ranks is equivalent, as an
unordered set of points, to selecting a uniform set of $k$ of those
independent points.  The selected unordered set therefore has the
law of $k$ independent uniform points on the circle.  Notice that the
fixed endpoint $t_0=0$ is not one of the selectable coordinates.

Partition the circle into $m'=\lfloor n/(2R)\rfloor$ equal cells.
Their length is at least $2R/n$, and $m'\ge n/(4R)$.  Let $p_a$ be
the number of selected points in cell $a$, with cell indices read
cyclically.  A point in cell $a$ can have counted neighbors only in
that cell or its two adjacent cells.  Consequently
\[
 \prod_{i\in I}N_i(I)^{1/4}
 \le\prod_{a:p_a>0}(p_a+p_{a-1}+p_{a+1})^{p_a/4}.
\]
The comparison with the product of the occupancies themselves follows
from
\[
 \begin{split}
 \sum_{a:p_a>0}p_a
  \log\frac{p_a+p_{a-1}+p_{a+1}}{p_a}
 &\le\sum_{a:p_a>0}(p_{a-1}+p_{a+1})\\
 &\le2k.
 \end{split}
\]
Thus the preceding product is at most
$e^{k/2}\prod_a p_a^{p_a/4}$, with $0^0=1$.

The occupancy vector has multinomial probabilities
$k!(m')^{-k}\prod_a(p_a!)^{-1}$.  The factorial lower bound implies
$p^{p/4}\le e^{p/4}(p!)^{1/4}$, so
\[
 \mathbb E\prod_a p_a^{p_a/4}
 \le e^{k/4}\frac{k!}{(m')^k}
    \sum_{\sum p_a=k}\left(\prod_a\frac1{p_a!}\right)^{3/4}.
\]
There are $\binom{k+m'-1}{k}$ nonnegative occupancy vectors.  Applying
H\"older's inequality to the last sum, and then the multinomial
theorem, gives
\[
 \sum_{\sum p_a=k}\left(\prod_a\frac1{p_a!}\right)^{3/4}
 \le
 \left(\frac{(m')^k}{k!}\right)^{3/4}
 \binom{k+m'-1}{k}^{1/4}.
\]
After substitution, the remaining factor apart from $C^k$ is
\[
 \left[\frac{k!}{(m')^k}\binom{k+m'-1}{k}\right]^{1/4}
 =\left[\prod_{q=0}^{k-1}(1+q/m')\right]^{1/4}
 \le(1+k/m')^{k/4}.
\]
Finally, $k\le n$ and $m'\ge n/(4R)$ give $k/m'\le4R$.
Since $R\ge1$, all factors are bounded as in
\eqref{pf:count-average}.
\end{proof}

Combining \eqref{pf:raw} with Lemma~\ref{pf:occupancy} gives a factor
of size $(CR^{-1/4})^k$, once the length factors have been extracted
and the coordinate subset has been averaged.  The remaining task is
to combine this subset average with the pointwise estimates for the
smooth stages.  We first make their scale dependence uniform.

\subsection{Expansion of the normalized curvature}

First record the dependence of the preceding estimates on $R$.
Since $r$ is a fixed multiple of $R$, the stopping condition compares
\[
 \log\beta_j=(3/2)^{j-1}\log R,
 \qquad
 \log(W^2R^{100})=O(R+\log R).
\]
Taking $j$ of order $1+\log R$ makes the first expression at least
the second.  As in \eqref{co:stage-count}, minimality therefore gives
\begin{equation}
 b_*=O(1+\log R).
 \label{pf:stage-count}
\end{equation}
For $j<J_*$, the constants $a_j$ are bounded independently of $R$,
because this is a fixed finite set of stages.  Their bases in
\eqref{pf:intermediate} are at most $CR^{-1/8}$.  At the later
stages,
\[
 a_j\le CR^{10}(1+b_*)^{20},
 \qquad \beta_j\ge R^{200},
\]
so the intermediate bases are at most
$CR^{-15}(1+\log R)^{20}$.  At the last stage, the stopping condition
$\beta_{b_*}\ge W^2R^{100}$ gives
\[
 Ca_{b_*}\sqrt W\,\beta_{b_*}^{-1/2}
 \le CR^{-40}(1+\log R)^{20}W^{-1/2}.
\]
Together with the raw factor just obtained, all these bases are at
most
\begin{equation}
 \varepsilon_R=C_1R^{-1/10}\le1
 \label{pf:small-base}
\end{equation}
for sufficiently large $R$, after choosing a fixed $C_1$ large enough.

The normalization of the connection gives
\begin{equation}
 d\alpha=C_{\rm nor}^{-1}\sum_{j<b_*}dB_j
       +C_{\rm nor}^{-1}W\,dB_{b_*}+Q,
 \qquad
 Q=-C_{\rm nor}^{-2}dC_{\rm nor}\wedge B.
 \label{pf:normalized-curvature}
\end{equation}
The lower bound $C_{\rm nor}\ge c_0>0$ from
Proposition~\ref{co:connection} changes the bases in
\eqref{pf:small-base} only by a fixed factor; enlarge $C_1$
accordingly.  The last term in \eqref{pf:normalized-curvature} is
decomposable, so $Q\wedge Q=0$.  Moreover,
\eqref{co:coefficients} implies
\begin{equation}
 |Q_{ii'}|\le C_Rl_i l_{i'}.
 \label{pf:rank-two}
\end{equation}
In particular, this stage can contribute to a Pfaffian partition
only on zero or two indices.  The constant $C_R$ in
\eqref{pf:rank-two} is independent of $n$.

There are $S=b_*+2$ stages in \eqref{pf:normalized-curvature}:
the raw stage, the $b_*$ smooth stages, and $Q$.  Denote their matrices
by $M_0,\ldots,M_{S-1}$, with $M_{S-1}=Q$.  The Pfaffian expansion is
\begin{equation}
 \Pf\left(\sum_{j=0}^{S-1}M_j\right)
 =\sum_{I_0\sqcup\cdots\sqcup I_{S-1}=\{1,\ldots,n\}}
       \sigma(I_0,\ldots,I_{S-1})
       \prod_{j=0}^{S-1}\Pf_{I_j}(M_j),
 \label{pf:partition-expansion}
\end{equation}
where all $|I_j|$ are even and each sign $\sigma$ is $1$ or $-1$.
There is no multiplicity factor for a fixed ordered partition.  To
check this normalization, expand the divided exterior power
$(\sum_jM_j)^m/m!$ as
$\sum_{\sum m_j=m}\prod_j M_j^{m_j}/m_j!$, and then extract the
coordinate volume coefficient using \eqref{pf:normalization} on
each coordinate subset.

Take absolute values in \eqref{pf:partition-expansion}.  In each
term the length factors multiply to $\prod_{i=1}^nl_i$, independently
of the partition.  Fix now the even sizes $k_j=|I_j|$.  For any given
raw subset $I_0$ of size $k_0$, the number of completions to a partition
of these sizes is
\[
 \frac{(n-k_0)!}{\prod_{j=1}^{S-1}k_j!},
\]
which is independent of $I_0$.  Every factor other than the raw
neighbor counts has already been bounded independently of the
simplex point and of the actual subset.  Thus, upon summing over
these partitions and integrating over the simplex, the count product
in \eqref{pf:raw} is averaged uniformly over all raw subsets of
size $k_0$.  Lemma~\ref{pf:occupancy} applies without any assertion
about an individual fixed subset.

The smooth estimates also contain factors
$(n/k_j)^{k_j/4}$ or $(n/k_j)^{k_j/2}$.  Their product is at most
\begin{equation}
 \prod_{j:k_j>0}(n/k_j)^{k_j/2}\le S^{n/2}.
 \label{pf:entropy}
\end{equation}
For completeness, put $p_j=k_j/n$ on the positive sizes.  Concavity
of the logarithm gives
\[
 \sum_jp_j\log(1/p_j)
 \le\log\left(\sum_jp_j/p_j\right)\le\log S;
\]
exponentiating $n/2$ times this inequality proves
\eqref{pf:entropy}.  Factors for the raw and decomposable stages may
be inserted here because they are at least one.

Let $q=k_{S-1}$.  A nonzero term has $q=0$ or $q=2$.  The small bases
in the other stages give $\varepsilon_R^{n-q}$, while the
decomposable stage costs at most $C_R$.  Since
$0<\varepsilon_R\le1$, both possibilities are bounded by
$C_R\varepsilon_R^n$ after absorbing the single factor
$\varepsilon_R^{-2}$ into $C_R$.  This absorption introduces no
dependence on $n$.

For fixed sizes, the number of ordered partitions is
$n!/\prod_j k_j!$.  Summed over all admissible sizes this is at most
$S^n$, the number of all assignments of the $n$ indices to $S$
stages.  Combining this count, \eqref{pf:entropy}, and the preceding
uniform subset average proves
\[
 \mathbb E_{\Delta_n}|\Pf(d\alpha)|
 \le C_R\varepsilon_R^n S^{3n/2}\prod_{i=1}^n l_i.
\]
Substitution of \eqref{pf:small-base} proves
\eqref{pf:main-estimate}, and \eqref{pf:stage-count} supplies the
assertion about $S$.  This completes the proof of
Theorem~\ref{pf:bound}.
